\newcommand{\tipoRequerimento}{Cobrança indevida sobre enquadramento tarifário incorreto}
\section{DOS FATOS E DOS DIREITOS}

% Contestação geral para cargas presentes que divergem da exclusividade de circuitos de iluminação

\begin{enumerate}[resume]
    \item O indeferimento do enquadramento na Iluminação Pública fundamentou-se na alegação de existência de apenas fiação e tomadas, sem nenhum equipamento conectado ou potência identificável, reputada pela distribuidora como indício de carga não exclusiva.

    \item Todavia, não foi identificado qualquer equipamento elétrico efetivamente conectado ao referido ponto, tampouco indicada potência nominal associada a equipamento instalado que pudesse integrar o cálculo da carga instalada da unidade consumidora.

    \item Igualmente, não houve medição, aferição técnica ou demonstração objetiva de consumo vinculado ao ponto em questão, inexistindo registro de demanda ativa ou reativa que evidenciasse utilização autônoma de energia elétrica.

    \item A evidência apresentada restringe-se à constatação da existência de elemento estrutural da instalação elétrica (fiação/soquete/tomada), desacompanhado de carga ativa, isto é, sem qualquer equipamento acoplado em condição de funcionamento.

    \item Do ponto de vista técnico, trata-se de componente passivo da instalação elétrica, desprovido de potência nominal própria, incapaz de demandar energia elétrica sem a conexão de equipamento externo.

    \item Ressalte-se, ainda, que o referido soquete encontra-se fixado/pendurado no poste, ao lado da caixa de medição, sem proteção/carcaça e exposto às intempéries, circunstâncias que evidenciam a inviabilidade prática de sua utilização para conexão de qualquer equipamento elétrico de uso geral, restringindo-se, na prática, à instalação de lâmpadas (ou componentes diretamente vinculados ao sistema de iluminação pública).

\end{enumerate}

\section{ANÁLISE REGULATÓRIA}

\begin{enumerate}[resume]
    \item Nos termos do art. 2º, III, da REN nº 1.000/2021, carga instalada é soma das potências nominais dos equipamentos elétricos instalados na unidade consumidora e em condições de entrar em funcionamento.

    \Citacao{
        Art. 2º Para os fins e efeitos desta Resolução, são adotadas as seguintes definições:

        [...]

        III - carga instalada: soma das potências nominais dos equipamentos elétricos instalados na unidade consumidora e em condições de entrar em funcionamento, expressa em kW (quilowatts);
    }

    \item Diante disso, impõe-se registrar que a interpretação do conceito de carga instalada; seja pela distribuidora, seja no âmbito da própria ANEEL; deve necessariamente se ancorar no regramento expresso da REN ANEEL nº 1.000/2021, observando-se seus dispositivos aplicáveis e a finalidade regulatória por eles tutelada, vedadas ampliações interpretativas dissociadas do texto normativo.

    \item Da definição normativa constante do art. 2º, III, da REN nº 1.000/2021, extraem-se requisitos cumulativos e indissociáveis para a caracterização de carga instalada, quais sejam:

    \begin{enumerate}
        \item a existência de equipamento elétrico efetivamente instalado na unidade consumidora;

        \item a presença de potência nominal tecnicamente identificável, apta a ser quantificada em kW; e

        \item a condição objetiva de entrada em funcionamento, isto é, a aptidão concreta do equipamento para operar e demandar potência do sistema elétrico.
    \end{enumerate}

    \item A ausência de qualquer desses elementos descaracteriza, sob o prisma técnico-regulatório, o enquadramento como carga instalada.

    \item Componentes como fios, condutores, tomadas e soquetes configuram elementos passivos da instalação elétrica, destinados exclusivamente à condução, derivação e conexão da energia no interior da unidade consumidora.

    \item Tais componentes não se qualificam como equipamentos elétricos de utilização final; não possuem potência nominal própria associada a consumo; não realizam conversão, transformação ou aproveitamento de energia elétrica; não demandam potência ativa (kW) nem potência reativa (kVAr) do sistema, quando desacompanhados de carga conectada.

    \item Assim, sob a perspectiva técnico-elétrica e à luz do conceito normativo de "carga instalada" previsto na Resolução Normativa nº 1.000/2021 da Agência Nacional de Energia Elétrica, tais componentes não se enquadram como carga, tampouco podem ser considerados para fins de aferição de exclusividade da destinação do fornecimento.
\end{enumerate}

\textbf{Interpretação do art. 189 --- Exclusividade}

\begin{enumerate}[resume]

    \item O art. 189 da REN nº 1.000/2021 dispõe o conceito e o âmbito de incidência da classe Iluminação Pública:

    \Citacao{
        Art. 189. Deve ser classificada na classe iluminação pública a unidade consumidora destinada exclusivamente à prestação do serviço público de iluminação pública, de responsabilidade do poder público municipal ou distrital ou daquele que receba essa delegação, com o objetivo de iluminar: (grifou-se)

        I - vias públicas destinadas ao trânsito de pessoas ou veículos, tais como ruas, avenidas, logradouros, caminhos, passagens, passarelas, túneis, estradas e rodovias; e

        II - bens públicos destinados ao uso comum do povo, tais como abrigos de usuários de transportes coletivos, praças, parques e jardins, ainda que o uso esteja sujeito a condições estabelecidas pela administração, inclusive o cercamento, a restrição de horários e a cobrança.
    }

    \item A expressão "destinada exclusivamente", constante do art. 189 da REN nº 1.000/2021 da Agência Nacional de Energia Elétrica, deve ser interpretada mediante critérios técnico-funcionais e finalísticos, compatíveis com a sistemática regulatória do enquadramento tarifário.

    \item Sob essa perspectiva, a caracterização de carga exclusiva pressupõe a verificação cumulativa de:

    \begin{enumerate}
        \item funcionalidade operacional - existência de vinculação direta ou instrumental à atividade-fim do serviço Iluminação Pública;

        \item instrumentalidade técnica - toda a carga instalada é destinada exclusivamente à alimentação dos refletores da praça;

        \item carga de fato instalada - a suposição quanto à existência de eventuais cargas que possam vir a ser instaladas no futuro, e que eventualmente descaracterizariam a destinação para iluminação pública, não constitui critério válido para o levantamento e a análise das cargas efetivamente instaladas no momento da verificação.
    \end{enumerate}

    \item Conforme demonstrado na evidência apresentada pela própria distribuidora, a suposta carga instalada limita-se à existência de fiação e soquete no local, sem qualquer equipamento conectado ou carga efetivamente instalada. Não há identificação de potência, tampouco indício de consumo mensurável que indique utilização diversa da finalidade originalmente atribuída. Diante desse cenário, a constatação técnica aponta para a inexistência de comprovação de desvio de finalidade do fornecimento, uma vez que a mera presença de infraestrutura elétrica (tomada/soquete), desacompanhada de uso efetivo, não caracteriza irregularidade, tampouco carga instalada.

    \item Cumpre destacar, ainda, que a classificação da unidade consumidora deve se apoiar em dados objetivos e verificáveis, apurados no ato da ligação e/ou na data da vistoria técnica, com base em evidências concretas (medições, inspeção \textit{in loco} e registros formais). Não se admite, portanto, que o enquadramento seja definido ou alterado com fundamento em meras suposições, presunções genéricas ou "achismos", até porque inexiste previsão regulatória que autorize a adoção de critério especulativo para afastar a classificação adequada.
\end{enumerate}

\textbf{Critério técnico-jurídico para caracterização de "carga não exclusiva"}

\begin{enumerate}[resume]

    \item A configuração de carga não exclusiva, apta a afastar o enquadramento prevista no art. 189 da REN nº 1.000/2021 da Agência Nacional de Energia Elétrica, exige a demonstração objetiva, técnica e cumulativa dos seguintes elementos:

    \begin{itemize}
        \item existência de equipamento elétrico efetivamente instalado na unidade consumidora;

        \item potência nominal identificável, passível de quantificação e integração ao cálculo da carga instalada;

        \item condição concreta de funcionamento autônomo, com possibilidade real de operação independente da atividade-fim;

        \item destinação diversa ou estranha ao serviço de iluminação pública, evidenciando utilização não vinculada à finalidade exclusiva exigida pela norma.
    \end{itemize}

    \item A ausência de qualquer desses requisitos impede a caracterização de carga não exclusiva.

    \item A simples presença de elemento estrutural da instalação elétrica, como ponto de tomada, soquete ou derivação, desacompanhado de equipamento conectado e sem demonstração de potência demandada, não satisfaz os critérios técnicos necessários à descaracterização da exclusividade.

    \item Adotar entendimento diverso implicaria em ampliação indevida do conceito regulatório de carga instalada, na criação de requisito não previsto expressamente na REN nº 1.000/2021, em utilização de presunção desprovida de base técnico-fática e em mitigação da segurança jurídica e da previsibilidade na aplicação da classificação tarifária.

\end{enumerate}

\section{DOS PEDIDOS}

\begin{enumerate}[resume]
\item 
Diante do exposto, requer que a Concessionária:

    \begin{enumerate}
        \item Devolva os valores cobrados indevidamente, referentes ao enquadramento da Unidade Consumidora.

        \item Proceda com a devolução de todos os valores cobrados a maior nos últimos 10 (dez) anos, conforme o Despacho nº 2.006 da Aneel, de 08 de julho de 2024, observando o prazo prescricional de 10 anos, conforme estabelecido no art. 205 do Código Civil. Conforme fundamentado na sentença proferida em 29 de setembro de 2023, pela 19ª Vara Cível Federal de São Paulo, no âmbito da Ação Civil Pública Cível nº 5024153-93.2018.4.03.6100\textsuperscript{1}, no Ofício nº 00710/2024/PFANEEL/PGF/AGU\textsuperscript{2}, de 08 de maio de 2024 e no Ofício nº 00937/2024/PFANEEL/PGF/AGU, de 12 de junho de 2024\textsuperscript{3}.

        \item Que a devolução se dê por depósito em conta corrente de titularidade do Município, na forma do art. 323, § 7º, II, da REN nº 1.000/2021, comunicando-se o cumprimento por escrito ao representante que a esta subscreve;

        \item Que todas as comunicações, intimações e notificações relativas à presente reclamação sejam encaminhadas exclusivamente ao representante legal do Município que a subscreve, preferencialmente por meio do endereço eletrônico \emailEmpresa{}, em observância aos princípios do contraditório, ampla defesa, eficiência e celeridade, bem como ao art. 9º da REN ANEEL nº 1.000/2021.
    \end{enumerate}

\end{enumerate}